\documentclass[11pt]{amsart}
\usepackage{amssymb}
\usepackage{booktabs}
\usepackage[expansion=false]{microtype}
\usepackage{xcolor}
\usepackage{xurl}
\usepackage[colorlinks=true,linkcolor=blue!50!black,citecolor=blue!50!black,urlcolor=blue!50!black]{hyperref}

\newtheorem{theorem}{Theorem}[section]
\newtheorem{lemma}[theorem]{Lemma}
\newtheorem{proposition}[theorem]{Proposition}
\newtheorem{corollary}[theorem]{Corollary}
\newtheorem{conjecture}[theorem]{Conjecture}
\newtheorem*{theoremA}{Theorem A}
\newtheorem*{theoremB}{Theorem B}
\newtheorem*{theoremC}{Theorem C}
\theoremstyle{definition}
\newtheorem{definition}[theorem]{Definition}
\theoremstyle{remark}
\newtheorem{remark}[theorem]{Remark}
\newtheorem{example}[theorem]{Example}

\newcommand{\erfc}{\operatorname{erfc}}
\newcommand{\Li}{\operatorname{Li}}
\newcommand{\arctanh}{\operatorname{arctanh}}
\newcommand{\dd}{\,\mathrm{d}}
\newcommand{\QQ}{\mathbb{Q}}
\newcommand{\Qbar}{\overline{\mathbb{Q}}}
\newcommand{\ZZ}{\mathbb{Z}}
\newcommand{\CC}{\mathbb{C}}
\newcommand{\RR}{\mathbb{R}}
\newcommand{\Mul}{\mathsf{M}}
\newcommand{\CT}{\operatorname*{CT}}
\newcommand{\Res}{\operatorname*{Res}}
\newcommand{\todo}[1]{\textcolor{red}{[TODO: #1]}}

% Working title; nomenclature to be decided ("zeta function of a holonomic function",
% "holonomic zeta function", "Dirichlet series of a holonomic function", ...).
\title[Zeta functions of holonomic functions]{Zeta functions of holonomic functions}
\author{Fredrik Johansson}
\date{Draft, November 2026}

\raggedbottom
\begin{document}

\begin{abstract}
For a holonomic function $F$ that is regular at infinity we study the Dirichlet series
$\zeta_F(s) = \sum_{k=1}^{\infty} F(k)\,k^{-s}$. This class contains the Riemann zeta function, Dirichlet
$L$-functions, the Dirichlet series attached to rational functions, and generating functions for many
classical constants: Khinchin's constant is $\exp(-Z'(0))$ where $Z(s) = \sum_{k=1}^{\infty} p_k k^{-s}$ is the Mellin
transform of the Gauss--Kuzmin distribution, and the H\"older means of the partial quotients are the values
$Z(-p)^{1/p}$. We give an exact approximate functional equation for $\zeta_F$, obtained by generalizing
Riemann's theta-function splitting: both sides are rapidly convergent sums of values of holonomic functions,
and the formula gives the meromorphic continuation of $\zeta_F$ to $\CC$ with explicit principal parts. We
show that $\zeta_F$ grows like a shifted Riemann zeta function in vertical strips, satisfies no linear difference
or differential equation in $s$, has $(T/2\pi)\log T + O(T)$ zeros of height at most $T$ to the right of an
explicit line, and obeys a Hardy--Littlewood type approximate functional equation whose dual
coefficients are the values $F(t/2\pi m)$. Finally, when $F$ is defined over $\Qbar$, $\zeta_F$ and its
derivatives can be computed to $p$ bits at algebraic points in $O(p^{1/2}\Mul(p)\log^2 p)$ bit operations,
which gives subquadratic algorithms for Khinchin's constant and the other constants above.
\end{abstract}

\maketitle

%======================================================================
\section{Introduction}

\subsection{Two constants}
Euler's constant and Khinchin's constant have strikingly similar series,
\begin{align*}
  \gamma &= \sum_{k=1}^{\infty}\Bigl(\frac1k - \log\Bigl(1+\frac1k\Bigr)\Bigr), \\
  \log K\,\log 2 &= -\sum_{k=2}^{\infty}\log\Bigl(1-\frac1k\Bigr)\log\Bigl(1+\frac1k\Bigr),
\end{align*}
both of the form $\sum_{k=k_0}^{\infty} F(k)$ with $F$ holonomic and analytic at infinity. Euler's constant also has a
second life as a coefficient of a Dirichlet series: $\zeta(s) = 1/(s-1) + \gamma + O(s-1)$. The same is
true of Khinchin's constant. With the Gauss--Kuzmin probabilities
$p_k = \log_2(1 + 1/(k(k+2)))$ of the partial quotients of a random real number, let
\begin{equation}\label{eq:Z}
  Z(s) = \sum_{k=1}^{\infty} p_k\,k^{-s} = \mathbb{E}\bigl[a_1^{-s}\bigr],
\end{equation}
the Mellin transform of the distribution of a partial quotient $a_1$ under the Gauss measure. Then $Z(0) = 1$
(the product $\prod_k (k+1)^2/(k(k+2))$ telescopes to $2$), $K = \prod_k k^{p_k} = \exp(-Z'(0))$, and the
H\"older means of the partial quotients are $K_p = Z(-p)^{1/p}$ for $p < 1$, $p \ne 0$; in particular the
harmonic mean is $K_{-1} = 1/Z(1)$. The function $Z$ is meromorphic in $\CC$, with simple poles at
$s = -1, -2, -3, \dots$ (Section~\ref{sec:khinchinzeta}).

More generally, for a sequence $F(k)$ that agrees for large $k$ with a holonomic function regular at infinity
(Definition~\ref{def:regular}) we consider
\begin{equation}\label{eq:zetaF}
  \zeta_F(s) = \sum_{k=1}^{\infty} F(k)\,k^{-s}.
\end{equation}
For $F = 1$ this is the Riemann zeta function; periodic $F$ give Dirichlet $L$-functions; rational $F$ give the
Dirichlet series attached to rational functions studied since Mellin and Mahler \cite{Mah28,CN81,Sar84,Pet98,Ess97}.
A sum $S = \sum_{k=1}^{\infty} F(k)$ is the value $\zeta_F(0)$, a sum $\sum_{k=1}^{\infty} (\log k)^j F(k)$ is
$(-1)^j\zeta_F^{(j)}(0)$, a divergent sum $\sum_{k=1}^{\infty} k^m F(k)$ has the zeta-regularized value $\zeta_F(-m)$, and
products $\prod_k k^{F(k)}$ are $\exp(-\zeta_F'(0))$. All constants in Table~\ref{tab:constants} are of these
forms.

The point of view of this paper is that $\zeta_F$, as a function of $s$, is a natural object that behaves in
many ways like a (shifted) Riemann zeta function, and that the classical theta-function method for $\zeta$
extends to it essentially unchanged, except that the functional equation is lost.

\subsection{Results}
Let $F$ be analytic at infinity, $F(k) = f(1/k)$ for $k \ge k_0$, where
$f(x) = \sum_{n=0}^{\infty} a_nx^n$ has radius of convergence $\rho > 1/k_0$.

\begin{theoremA}[exact approximate functional equation; Theorem~\ref{thm:split}]
For every $u_0 > 0$ and every $s \in \CC$ except the poles,
\[
  \zeta_F(s) = \sum_{k=1}^{k_0-1}F(k)k^{-s} + \sum_{k=k_0}^{\infty}\mathcal{L}_s(k^2)
  + \sqrt\pi\sum_{m=1}^{\infty}\mathcal{D}_s(\pi^2m^2) + M(s),
\]
where $\mathcal{L}_s(K) = \int_{u_0}^\infty e^{-Ku}h_s(u)\dd u$ and
$\mathcal{D}_s(c) = \int_0^{u_0}u^{-1/2}e^{-c/u}h_s(u)\dd u$ are entire in $s$ and holonomic in $K$ and $c$,
$h_s(u) = \sum_{n=0}^{\infty} a_n u^{(s+n)/2-1}/\Gamma(\frac{s+n}{2})$ is a Borel-type transform of $f$, and $M$ is an
explicit series of elementary functions and incomplete gamma functions carrying all the poles.
\end{theoremA}

For $f = 1$, $k_0 = 1$ and $u_0 = \pi$ this is Riemann's formula, and both sides are the same function: the
functional equation. In general the two sides differ, but both are rapidly convergent: the $k$-terms decay
like $e^{-k^2u_0}$ and the dual terms like $e^{-\pi^2m^2/u_0}$. Formally the dual side is the Poisson dual
$\sum_{m\in\ZZ} \widehat G_s(m)$ of $G_s(x) = |x|^{-s}F(|x|)$, and $u_0$ is a smoothing parameter
(Remark~\ref{rem:poisson}).

\begin{theoremB}[Section~\ref{sec:properties}, Section~\ref{sec:larget}]
Let $n_0$ be the least $n$ with $a_n \ne 0$.
\begin{enumerate}
\item $\zeta_F$ is meromorphic in $\CC$ with at most simple poles at $s = 1-n$, of residue $a_n$, and
$\zeta_F(s) - \int_{k_0}^\infty F(x)x^{-s}\dd x$ is entire.
\item $\zeta_F$ satisfies no linear difference or differential equation in $s$ with polynomial coefficients
(unless $F$ has finite support).
\item The growth exponent of $\zeta_F$ on vertical lines is $\mu_F(\sigma) = \mu(\sigma+n_0)$, where $\mu$ is that
of $\zeta$. In particular the Lindel\"of hypothesis is equivalent to $\zeta_F(\frac12 - n_0 + it) \ll t^\varepsilon$.
\item For every $\varepsilon > 0$, $\zeta_F$ has $(T/2\pi)\log(T/2\pi e k_1) + O(\log T)$ zeros with
$0 < \operatorname{Im}s \le T$ and $\operatorname{Re}s \ge -n_0-\varepsilon$, and none with
$\operatorname{Re}s \le -n_0-\varepsilon$ and large imaginary part in any fixed strip; here $k_1$ is the least $k$
with $F(k) \ne 0$.
\end{enumerate}
\end{theoremB}

Numerically, the zeros of Khinchin's $Z$ accompany the zeros of $\zeta(s+2)$, displaced to the right by an amount
that decreases slowly with the height (Section~\ref{sec:zeros}); we conjecture that almost all of them lie
arbitrarily close to the line $\operatorname{Re}s = \frac12 - n_0$, as Levinson proved for the $a$-points of $\zeta$
\cite{Lev75}. We also give a Hardy--Littlewood approximate functional equation
\[
  \zeta_F(s) \approx \sum_{1\le k\le x}F(k)k^{-s} + \chi(s)\sum_{1\le m\le y}F\Bigl(\frac{t}{2\pi m}\Bigr)m^{s-1},
  \qquad 2\pi xy = t,
\]
in which the dual coefficients are the values of $F$ at the stationary points (Proposition~\ref{prop:hl}).

\begin{theoremC}[Theorem~\ref{thm:complexity}]
If $F$ is regular at infinity and defined over $\Qbar$, then for every algebraic $s_0$ and fixed $j$ the
$j$-th Taylor (or Laurent) coefficient of $\zeta_F$ at $s_0$ can be computed to $p$ bits in
$O(p^{1/2}\Mul(p)\log^2 p)$ bit operations.
\end{theoremC}

Here $\Mul(p)$ is the cost of $p$-bit integer multiplication, so the bound is $\widetilde O(p^{3/2})$, against
$p^{2+o(1)}$ for the classical methods (Euler--Maclaurin summation, the expansion in zeta values). The
algorithmic part, and computations of Khinchin's constant to $10^7$ digits and of the other constants of
Table~\ref{tab:constants} to $10^6$ digits, are described in the companion paper \cite{Joh26}; here we
concentrate on $\zeta_F$ as a function.

\subsection{Related work}
The meromorphic continuation of a Dirichlet series $\sum_{k=1}^{\infty} F(k)k^{-s}$ whose coefficients have an asymptotic
expansion at infinity is classical; it follows from Euler--Maclaurin summation or the Mellin transform method
\cite{Zag06,FGD95}, and is used routinely for zeta regularization in physics \cite{Eli95,Kir02}. Dirichlet series
attached to polynomials and rational functions, in one and several variables, have been studied by Mahler
\cite{Mah28}, Cassou-Nogu\`es \cite{CN81}, Sargos \cite{Sar84}, Peter \cite{Pet98}, Essouabri \cite{Ess97}
and others; Apostol and Vu \cite{AV84} treat series such as $\sum_{k=1}^{\infty} k^{-s}(k+1)^{-1}$. Mellin--Barnes integrals
give the continuation of hypergeometric cases such as $\sum_{k=1}^{\infty} k^{-s}(k+a)^{-t}$ and their asymptotic
expansions in parameters \cite{Kat97}. Approximate functional equations go back to Riemann \cite{Rie59} and
Hardy and Littlewood \cite{HL23}; the theta-function (``smoothed'') form for general $L$-functions is due to
Lavrik \cite{Lav68} and is the basis of the algorithms of Rubinstein \cite{Rub05} and Dokchitser \cite{Dok04};
for Dirichlet $L$-functions at algebraic points it gives $p^{3/2+o(1)}$ algorithms \cite{Joh21}, which
Theorem~C generalizes.
Crandall \cite{Cra12} developed the splitting systematically for polylogarithms, Lerch transcendents and
related functions. Zeros of Dirichlet series without Euler product, such as the Davenport--Heilbronn examples
\cite{DH36}, and value distribution in general, are surveyed in \cite{Ste07}. To our knowledge the class of
holonomic $F$ has not been studied as such; what holonomy adds is that every piece of the splitting is again
holonomic, which is what makes the fast evaluation possible.

\subsection{Outline}
Section~\ref{sec:defs} defines the regular class and $\zeta_F$, and lists examples. Section~\ref{sec:properties}
proves the elementary properties: continuation and principal parts, the Mellin part, values at
non-positive integers, growth and mean values, and the absence of linear difference and differential
equations. Section~\ref{sec:splitting} proves Theorem~A. Section~\ref{sec:larget} is about large $|\operatorname{Im}s|$:
the Hardy--Littlewood approximate functional equation, zero-free regions, zero counting, and numerical
experiments with Khinchin's $Z$. Section~\ref{sec:complexity} proves Theorem~C. Section~\ref{sec:special} discusses
the special values in Table~\ref{tab:constants}, the Khinchin zeta function and a two-variable family that
contains both $\zeta$ and $Z$, and Section~\ref{sec:open} lists open problems.

%======================================================================
\section{Definitions and examples}\label{sec:defs}

\subsection{The regular class}
Let $L = \sum_{i=0}^r p_i(x)\,\partial_x^i$ with $p_i \in \CC[x]$, $p_r \ne 0$, and suppose that $x = \infty$ is
an ordinary point or a regular singular point of $L$, i.e.\ that $t = 1/x$ is such a point of the transformed
operator. Let $R$ exceed the absolute values of the finite singular points of $L$. By Frobenius's method
\cite[\S16]{Inc26}, every solution of $L$ on $(R, \infty)$ can be written as
\begin{equation}\label{eq:frob}
  F(x) = \sum_{\sigma\in\Sigma}\sum_{j=0}^{J} x^{-\sigma}(\log x)^j f_{\sigma,j}(1/x),
\end{equation}
with a finite set $\Sigma$ of exponents, pairwise incongruent modulo $\ZZ$, and power series $f_{\sigma,j}$
convergent for $|t| < 1/R$.

\begin{definition}\label{def:regular}
A function $F$ on $(R,\infty)$ is \emph{regular at infinity} if it is annihilated by an operator $L$ as above.
It is \emph{analytic at infinity} if moreover $F(x) = f(1/x)$ with $f$ analytic at $0$. It is \emph{defined over
$\Qbar$} if $L$ can be chosen with $p_i \in \Qbar[x]$ and the coefficients of the power series $f_{\sigma,j}$ in
\eqref{eq:frob} are algebraic numbers. A sequence $(F(k))_{k\ge1}$ is \emph{regular} if it agrees for $k \ge k_0$
with a function regular at infinity, for some integer $k_0 > R$.
\end{definition}

Regularity at infinity is preserved by sums, products, $F(x) \mapsto F(qx+a)$, multiplication by $x^{-\sigma}$ and
by $(\log x)^j$, differentiation and integration. It holds for all G-functions of $1/x$ (algebraic functions,
logarithms, polylogarithms, ${}_{q+1}F_q$ with rational parameters) and for entire functions of $1/x$ such as
$e^{c/x^2}$, whose irregular singularity is at $x = 0$ and is never reached. What it excludes is exponential
behaviour at infinity, such as $z^k$ with $|z| \ne 1$ or $e^{ik^2}$.

\begin{definition}
For a regular sequence $F$, let
\[
  \zeta_F(s) = \sum_{k=1}^{\infty}F(k)\,k^{-s}, \qquad
  \zeta_F(s;k_0) = \sum_{k=k_0}^{\infty}F(k)\,k^{-s},
\]
and, for a periodic $\chi\colon\ZZ \to \CC$, $L_F(s,\chi) = \sum_{k=1}^\infty \chi(k)F(k)k^{-s}$.
\end{definition}

The series converge absolutely for $\operatorname{Re}s > 1 - \beta$, where $\beta = \min\operatorname{Re}\sigma$ over
the exponents that occur. The first $k_0 - 1$ terms are entire functions of $s$, so all analytic questions concern
the tail $\zeta_F(s;k_0)$. (The notation $\zeta_F(s;k_0)$ shifts the index, not the argument; it is not a Hurwitz
zeta function.)

\begin{lemma}\label{lem:reduction}
If $F$ is given for $x \ge k_0$ by \eqref{eq:frob}, and $F_{\sigma,j}(x) = f_{\sigma,j}(1/x)$, then
\[
  \zeta_F(s;k_0) = \sum_{\sigma\in\Sigma}\sum_{j=0}^{J}(-\partial_s)^j\,\zeta_{F_{\sigma,j}}(s+\sigma;k_0).
\]
\end{lemma}

\begin{proof}
Termwise, since $x^{-\sigma}(\log x)^jx^{-s} = (-\partial_s)^jx^{-s-\sigma}$, for $\operatorname{Re}s$ large; then by
continuation.
\end{proof}

So everything reduces to $F$ analytic at infinity, at the price of shifts and derivatives in $s$, and from now
on we assume, unless stated otherwise, that
\begin{equation}\label{eq:f}
  F(k) = f(1/k) \quad (k \ge k_0), \qquad f(x) = \sum_{n=0}^{\infty}a_nx^n, \qquad |a_n| \le C R^n,
  \quad R < k_0,
\end{equation}
and we let $n_0 = \min\{n : a_n \ne 0\}$ (we exclude $f = 0$).

\subsection{Examples}
\begin{enumerate}
\item $F = 1$ gives $\zeta(s)$, $F = \chi$ periodic gives $L(s,\chi)$, and $F(k) = k^{-a}$ gives $\zeta(s+a)$.
\item $F$ rational gives the Dirichlet series attached to rational functions; for instance
$\sum_{k=2}^{\infty}k^{-s}/(k^2-1)$ and $\sum_{k=1}^{\infty}k^{-s}/(k(k+1))$, whose derivatives at $0$ are the
constants $\mathcal A$ and $\mathcal C$ of Table~\ref{tab:constants}.
\item The Khinchin zeta function \eqref{eq:Z}: $F(k) = p_k$, $f(x) = \log_2\bigl((1+x)^2/(1+2x)\bigr)$, $\rho = \frac12$.
\item Families: if $G$ is regular and $G(x) \to 1$ as $x \to \infty$, then $F_t(k) = G(k)^t$ is regular for every
$t$, which gives two-variable functions
\[
  \Phi(s,t) = \sum_{k=1}^{\infty} G(k)^tk^{-s}.
\]
The Hurwitz zeta function $\sum_{k=1}^{\infty}(k+a)^{-s} = \Phi(s,-s)$ with $G(x) = 1 + a/x$, and the Dirichlet
series $\sum_{k=1}^{\infty} P(k)^{-s}$ of a polynomial, are diagonals of such families. (The summand depends on $s$,
so they are not themselves of the form \eqref{eq:zetaF}.)
\item Not regular: the polylogarithm $\sum_{k=1}^{\infty} z^kk^{-s}$ with $|z|<1$, and the Lerch-type series
$\sum_{k=1}^{\infty} e^{2\pi i\alpha k}k^{-s}$ with irrational $\alpha$; the latter can be treated as a twist
(Remark~\ref{rem:twist}).
\end{enumerate}

\begin{table}[t]
\centering\footnotesize
\renewcommand{\arraystretch}{1.3}
\begin{tabular}{lll}
\toprule
constant & $F(k)$ & expression \\
\midrule
Euler's $\gamma$ & $1$; or $F_1(k)$ & $\CT_{s=1}\zeta(s)$; or $\zeta_{F_1}(0)$ \\
Stieltjes $\gamma_1$ & $F_1(k)$, $F_2(k)$ & $-\zeta_{F_1}'(0) - \frac12\zeta_{F_2}(0)$ \\
Khinchin's $K$ & $p_k = \log_2(1+\frac{1}{k(k+2)})$ & $\exp(-Z'(0))$ \\
H\"older means $K_p$, $p<1$ & $p_k$ & $Z(-p)^{1/p}$ \\
$\mathcal A = \sum_{k=2}^{\infty}\frac{\log k}{k^2-1}$ & $\frac1{k^2-1}$ ($k\ge2$) & $-\zeta_F'(0)$ \\
$\mathcal C = \sum_{k=1}^{\infty}\frac{\log k}{k(k+1)}$ & $\frac{1}{k(k+1)}$ & $-\zeta_F'(0)$ \\
vortex rings $U(\lambda)$ & ${}_2F_1(\frac32,\frac32;2;-\frac{4}{\lambda^2k^2})$ & $\lambda^{-3}\zeta_F(3)$ \\
Hardy--Littlewood $H(x)$ & $\sin(x/k)$ & $\zeta_F(1)$ \\
Planck--Larkin $\sigma(T)$ & $e^{x/k^2} - 1 - x/k^2$ & $\zeta_F(-2)$ \\
Casimir-type sums & $F(k)$ & $\zeta_F(-m)$ (regularized) \\
\bottomrule
\end{tabular}
\caption{Constants as special values, with $F_1(k) = \frac1k - \log(1+\frac1k)$ and
$F_2(k) = \log^2(1+\frac1k)$. The Alladi--Grinstead and Kalpazidou constants are $e^{\mathcal C-1}$ and
$e^{\mathcal C}$. For the vortex sum $\lambda$ is fixed and $k_0 > 2/\lambda$; for Planck--Larkin
$x = I/(k_BT)$.}
\label{tab:constants}
\end{table}

%======================================================================
\section{Elementary properties}\label{sec:properties}

\subsection{Continuation and principal parts}

\begin{proposition}\label{prop:cont}
Under \eqref{eq:f},
\begin{equation}\label{eq:zetavalue}
  \zeta_F(s;k_0) = \sum_{n=0}^{\infty}a_n\,\zeta(s+n,k_0), \qquad \zeta(w,k_0) = \sum_{k=k_0}^{\infty}k^{-w},
\end{equation}
where the series converges locally uniformly on $\CC \setminus \{1, 0, -1, \dots\}$. Hence $\zeta_F$ is meromorphic
in $\CC$, with at most simple poles at $s = 1-n$, $n \ge n_0$, of residue $a_n$. For a general regular $F$
\eqref{eq:frob}, the poles lie at $s = 1 - \sigma - n$ and have order at most $J+1$, with principal parts given
by Lemma~\ref{lem:reduction}.
\end{proposition}

\begin{proof}
For $\operatorname{Re}w = v > 1$, $|\zeta(w,k_0)| \le k_0^{-v}(1 + k_0/(v-1))$, so for $s$ in a compact set and
$n$ large the terms are $O((R/k_0)^n)$; for the finitely many remaining $n$, $\zeta(s+n,k_0)$ is meromorphic with a
single simple pole at $s = 1-n$ of residue $1$. For $\operatorname{Re}s > 1$ the double series converges absolutely,
which justifies the interchange.
\end{proof}

Formula \eqref{eq:zetavalue} is the \emph{zeta-value expansion}. It is the classical way to compute such sums,
but it needs $\Theta(p)$ zeta values to $p$ bits, hence $p^{2+o(1)}$ bit operations; the splitting of
Section~\ref{sec:splitting} replaces it.

\subsection{The Mellin part}
Let $\mathcal M_F(s;k_0) = \int_{k_0}^\infty F(x)x^{-s}\dd x$, the Mellin transform of $F$ restricted to $[k_0,\infty)$.

\begin{proposition}\label{prop:mellin}
Under \eqref{eq:f}, $\mathcal M_F(s;k_0) = \sum_{n=0}^{\infty} a_n k_0^{1-s-n}/(s+n-1)$, and $\zeta_F(s;k_0) - \mathcal M_F(s;k_0)$ is entire.
If $F$ is annihilated by $\sum_{i=-r}^{r} x^ir_i(\vartheta_x)$, $\vartheta_x = x\frac{\mathrm{d}}{\mathrm{d}x}$, then $\mathcal M_F$
satisfies the inhomogeneous linear recurrence
\[
  \sum_{i=-r}^{r} r_i(s-i-1)\,\mathcal M_F(s-i;k_0) = k_0^{1-s}\,P(s),
\]
where $P$ is a polynomial whose coefficients are linear in $F(k_0), F'(k_0), \dots$.
\end{proposition}

\begin{proof}
The first formula is termwise integration. For the second, apply the Euler--Maclaurin formula of order $2N$ to
$g(x) = F(x)x^{-s}$ on $[k_0,\infty)$: the difference is $\frac12g(k_0)$, plus finitely many derivatives of $g$
at $k_0$, plus $\int_{k_0}^\infty \widetilde B_{2N}(x)g^{(2N)}(x)\dd x/(2N)!$, and $g^{(2N)}(x) \ll x^{-\operatorname{Re}s-n_0-2N}$
locally uniformly in $s$, so the difference is analytic for $\operatorname{Re}s > 1 - n_0 - 2N$, for every $N$.
For the recurrence, $\int_{k_0}^\infty x^{-s}\vartheta_xF\dd x = -k_0^{1-s}F(k_0) + (s-1)\mathcal M_F(s;k_0)$ and multiplication by
$x^i$ shifts $s$ by $-i$; iterate.
\end{proof}

So all the poles of $\zeta_F$, and all of its ``holonomic'' structure in $s$, sit in the Mellin part. The
remainder is entire but, in contrast to $\mathcal M_F$, not holonomic in $s$:

\begin{proposition}\label{prop:nonholonomic}
Let $F$ be any sequence with $F(k) \ne 0$ for infinitely many $k$ and $\sum_{k=1}^{\infty}F(k)k^{-s}$ convergent somewhere.
Then $\zeta_F$ satisfies no relation $\sum_{i=0}^rp_i(s)\zeta_F(s+i) = q(s)$ and no relation
$\sum_{i=0}^rp_i(s)\zeta_F^{(i)}(s) = q(s)$ with polynomials $p_i$ not all zero and $q$ a finite
linear combination of exponentials $c^{-s}$ with polynomial coefficients.
\end{proposition}

\begin{proof}
For $\operatorname{Re}s$ large the left-hand sides are generalized Dirichlet series
\[
  \sum_{k=1}^{\infty} F(k)\,c_k(s)\,k^{-s}, \quad c_k(s) = \sum_{i=0}^{r}p_i(s)k^{-i}
  \ \text{or}\ c_k(s) = \sum_{i=0}^{r}p_i(s)(-\log k)^i,
\]
polynomial in $s$ of bounded degree. By the uniqueness theorem for such series (divide by the highest power of
$s$ and let $\operatorname{Re}s\to\infty$, inductively), $F(k)c_k(s) = 0$ for every $k$ not among the finitely many
frequencies of $q$. Hence $c_k(s) = 0$ for infinitely many $k$, and as a polynomial in $k^{-1}$ (respectively
$\log k$) with infinitely many roots, $\sum_{i=0}^{r}p_i(s)y^i \equiv 0$, so all $p_i = 0$.
\end{proof}

Whether $\zeta_F$ can satisfy an algebraic differential equation is a different matter: $\zeta$ itself does not
(Hilbert, Ostrowski \cite{Ost20}), and we expect the same for every $\zeta_F$ that is not a finite combination of
exponentials, but we have no proof.

\subsection{Values at non-positive integers}
Let $B_n(x)$ denote the Bernoulli polynomials and $\psi$ the digamma function.

\begin{proposition}\label{prop:negative}
Under \eqref{eq:f}, for every integer $m \ge 0$,
\begin{multline*}
  \CT_{s=-m}\zeta_F(s;k_0) = -\sum_{n=0}^{m}a_n\frac{B_{m-n+1}(k_0)}{m-n+1} - a_{m+1}\psi(k_0) \\
  + \sum_{k=k_0}^{\infty}k^m\Bigl(f(1/k) - \sum_{n=0}^{m+1}a_nk^{-n}\Bigr),
\end{multline*}
where $\CT$ is the constant term of the Laurent expansion (the value, if $a_{m+1} = 0$).
\end{proposition}

\begin{proof}
In \eqref{eq:zetavalue}, $\zeta(-j, k_0) = -B_{j+1}(k_0)/(j+1)$ for $j \ge 0$,
$\zeta(1+\varepsilon,k_0) = 1/\varepsilon - \psi(k_0) + O(\varepsilon)$, and the terms with $n \ge m+2$ sum to the
absolutely convergent last sum.
\end{proof}

So the zeta-regularized value of a divergent sum $\sum_{k=1}^{\infty} k^mF(k)$ is a convergent sum of the same type plus
rational (if the $a_n$ are rational) Bernoulli corrections. For Khinchin's $Z$, Proposition~\ref{prop:negative}
at $m = 1$ gives $\CT_{s=-1}Z(s) = -1/\log2$ (Section~\ref{sec:khinchinzeta}).

\subsection{Growth and mean values}
Let $\mu_F(\sigma) = \inf\{\mu : \zeta_F(\sigma+it) \ll |t|^\mu \text{ for } |t| \ge 1\}$, and let $\mu = \mu_\zeta$.

\begin{proposition}\label{prop:growth}
Under \eqref{eq:f}, $\mu_F(\sigma) = \mu(\sigma + n_0)$ for every real $\sigma$.
\end{proposition}

\begin{proof}
Fix $\sigma$ and choose $N$ with $1 - N < \sigma - 1$. By \eqref{eq:zetavalue},
\[
  \zeta_F(s) = \sum_{k<k_0}F(k)k^{-s} + \sum_{n=n_0}^{N-1}a_n\zeta(s+n,k_0) + R_N(s),
\]
where $R_N(s) = \sum_{k=k_0}^{\infty}k^{-s}\bigl(f(1/k) - \sum_{n=0}^{N-1}a_nk^{-n}\bigr)$ is bounded on
$\operatorname{Re}s = \sigma$, as are the finite sums, and $\zeta(w,k_0)$ differs from $\zeta(w)$ by a
bounded finite sum. Since $\mu$ is nonincreasing, $\mu_F(\sigma) \le \mu(\sigma+n_0)$. Conversely
$a_{n_0}\zeta(s+n_0) = \zeta_F(s) - \sum_{n=n_0+1}^{N-1}a_n\zeta(s+n) + O(1)$. If $\mu(\sigma+n_0) = 0$ there is nothing to
prove ($\mu_F \ge 0$ everywhere: $\mu_F$ is convex, and $\mu_F = 0$ far to the right, where $\zeta_F$ is an
absolutely convergent Dirichlet series that is not $o(1)$). If $\mu(\sigma+n_0) > 0$, then
$\mu(\sigma+n_0+1) < \mu(\sigma+n_0)$: with $x = \sigma+n_0$, for $x \le 0$ because $\mu(x) = \frac12 - x$ while
$\mu(x+1) \le \max(-\frac x2, -\frac12 - x)$, and for $x > 0$ because $\mu(x+1) = 0$. Hence
$\mu(\sigma+n_0) \le \max(\mu_F(\sigma), \mu(\sigma+n_0+1))$ forces $\mu(\sigma+n_0) \le \mu_F(\sigma)$.
\end{proof}

\begin{corollary}\label{cor:lindelof}
\begin{enumerate}
\item The Lindel\"of hypothesis holds if and only if $\zeta_F(\frac12 - n_0 + it) \ll |t|^\varepsilon$ for every
$\varepsilon > 0$.
\item $\displaystyle\int_1^T\bigl|\zeta_F(\tfrac12 - n_0 + it)\bigr|^2\dd t = |a_{n_0}|^2\,T\log T + O\bigl(T(\log T)^{1/2}\bigr)$.
\end{enumerate}
\end{corollary}

\begin{proof}
On the line $\operatorname{Re}s = \frac12 - n_0$, $\zeta_F(s) = a_{n_0}\zeta(s+n_0) + O(1)$, since
$\zeta(s+n)$ for $n > n_0$ is evaluated at real part at least $\frac32$. (1) is immediate, and (2) follows from
$\int_1^T|\zeta(\frac12+it)|^2\dd t \sim T\log T$ and the Cauchy--Schwarz inequality for the cross term.
\end{proof}

%======================================================================
\section{Theta splitting}\label{sec:splitting}

\subsection{The kernel}
Under \eqref{eq:f} define, for $s \in \CC$ and $u > 0$,
\begin{equation}\label{eq:kernel}
  h_s(u) = \sum_{n=0}^{\infty} a_n\frac{u^{(s+n)/2-1}}{\Gamma(\frac{s+n}{2})} = u^{s/2-1}E_s(u^{1/2}), \qquad
  E_s(v) = \sum_{n=0}^{\infty}\frac{a_nv^n}{\Gamma(\frac{s+n}2)}.
\end{equation}
Since $1/\Gamma$ is entire and $|a_n| \le CR^n$, $E_s(v)$ is entire in $(s,v)$, and for $s$ in a compact set $\mathcal K$,
$|E_s(v)| \le C(1+|v|)^{c}e^{R^2|v|^2}$. If $\operatorname{Re}(s+n) > 0$ for all $n$ with $a_n \ne 0$, then for
$k > R$
\begin{equation}\label{eq:laplace}
  k^{-s}f(1/k) = \int_0^\infty e^{-k^2u}h_s(u)\dd u,
\end{equation}
by $k^{-w} = \Gamma(w/2)^{-1}\int_0^\infty u^{w/2-1}e^{-k^2u}\dd u$ and Tonelli's theorem. Writing $u = v^2$,
$h_s$ is $u^{s/2}$ times an even and an odd entire function of $v$:
\begin{equation}\label{eq:kernelparts}
\begin{gathered}
  h_s(u) = u^{s/2}\bigl(g_0(u) + u^{-1/2}g_1(u)\bigr), \\
  g_0(u) = \sum_{m=0}^{\infty}\frac{a_{2m}u^{m-1}}{\Gamma(m+\frac s2)}, \qquad
  g_1(u) = \sum_{m=0}^{\infty}\frac{a_{2m+1}u^{m}}{\Gamma(m+\frac{s+1}2)}.
\end{gathered}
\end{equation}

\begin{lemma}[{\cite[Lemma~3.2]{Joh26}}]\label{lem:kernelhol}
If $f$ is holonomic and its operator is regular singular (or ordinary) at $x = 0$, then for every fixed $s$
the functions $u\,g_0(u)$ and $g_1(u)$ are holonomic, annihilated by operators with coefficients in
$\QQ(s)[u]$ (if $f$ has rational coefficients), and the vector
$G = (g, \vartheta g, \dots, \vartheta^{r-1}g)^T$ satisfies $\vartheta G = A(u)G + b(u)$ with polynomial
$A$ and $b$.
\end{lemma}

In the companion paper the lemma is stated for $s = 0$ and, in the proof of Theorem~3.12 there, for any fixed
$s$; the proof only uses that $\Gamma(m+\sigma_0+1)/\Gamma(m+\sigma_0)$ is rational in $m$ and $\sigma_0$.

\subsection{The splitting}
Let $Q(a,x) = \Gamma(a,x)/\Gamma(a)$ and $P(a,x) = 1 - Q(a,x) = \gamma(a,x)/\Gamma(a)$ be the regularized
incomplete gamma functions; for fixed $x > 0$ both are entire in $a$.

\begin{theorem}\label{thm:split}
Assume \eqref{eq:f} and let $u_0 > 0$. For $s \in \CC$ not a pole,
\[
  \zeta_F(s;k_0) = \sum_{k=k_0}^{\infty}I_k(s) + \sqrt\pi\sum_{m=1}^{\infty}J_m(s) + M(s),
\]
where
\begin{align*}
  I_k(s) &= \int_{u_0}^\infty e^{-k^2u}h_s(u)\dd u = \sum_{n=0}^{\infty}a_n\,k^{-s-n}\,Q\Bigl(\frac{s+n}2, k^2u_0\Bigr), \\
  J_m(s) &= \int_0^{u_0}u^{-1/2}e^{-\pi^2m^2/u}h_s(u)\dd u, \\
  M(s) &= \sum_{n=0}^{\infty}a_n\biggl(\frac{\sqrt\pi\,u_0^{(s+n-1)/2}}{(s+n-1)\,\Gamma(\frac{s+n}{2})}
     - \frac{u_0^{(s+n)/2}}{2\,\Gamma(\frac{s+n}2+1)} \\
     &\qquad\qquad - \sum_{k=1}^{k_0-1}k^{-s-n}P\Bigl(\frac{s+n}2, k^2u_0\Bigr)\biggr).
\end{align*}
The functions $I_k$ and $J_m$ are entire, and for $s$ in a compact set $\mathcal K$,
$I_k(s) \ll e^{-(k^2-R^2)u_0}$ and $J_m(s) \ll e^{-\pi^2m^2/u_0}$ uniformly; all series converge locally uniformly,
and the only singularities of $M$ are the simple poles of its first term at $s = 1-n$, with residue $a_n$.
\end{theorem}

\begin{proof}
First let $\operatorname{Re}s > 1$. By \eqref{eq:laplace}, $\zeta_F(s;k_0) = \int_0^\infty h_s(u)\Theta(u)\dd u$
with $\Theta(u) = \sum_{k=k_0}^{\infty}e^{-k^2u}$, the interchange being justified by Tonelli on $[u_0,\infty)$ and,
on $(0,u_0]$, by $\Theta(u) \le \sqrt{\pi/u}$ and $|h_s(u)| \ll u^{\operatorname{Re}s/2-1}$. Split the integral at $u_0$.
On $[u_0,\infty)$ integrate termwise in $k$, which gives the $I_k$, and termwise in $n$, which gives the series for
$I_k$ (with $\int_{u_0}^\infty e^{-Ku}u^{w-1}\dd u = K^{-w}\Gamma(w)Q(w,Ku_0)$). On $(0,u_0]$ use the theta
transformation
\[
  \Theta(u) = \frac12\sqrt{\frac\pi u} - \frac12 - \sum_{k=1}^{k_0-1}e^{-k^2u} + \sqrt{\frac\pi u}\sum_{m=1}^{\infty}e^{-\pi^2m^2/u}.
\]
The last sum gives the $J_m$. The first three terms give $M$, by termwise integration in $n$ of
$\int_0^{u_0}u^{w-3/2}\dd u = u_0^{w-1/2}/(w-\frac12)$, $\int_0^{u_0}u^{w-1}\dd u = u_0^w/w$ and
$\int_0^{u_0}e^{-Ku}u^{w-1}\dd u = K^{-w}\Gamma(w)P(w,Ku_0)$ with $w = (s+n)/2$.

Now the continuation. $I_k$ is entire since $u \ge u_0$, and $|h_s(u)| \ll e^{R^2u}$ there gives the bound.
$J_m$ is entire since $e^{-\pi^2m^2/u}$ kills any power of $u$ at $0$, and
$e^{-\pi^2m^2/u} \le e^{-\pi^2(m^2-1)/u_0}e^{-\pi^2/u}$ gives the bound. In $M$, $\Gamma(w)^{-1}$ cancels the
pole of $u_0^w/w$, and $k^{-s-n}P(\frac{s+n}2,k^2u_0) \ll (k^2u_0)^{n/2}k^{-n}/\Gamma(\frac n2) = u_0^{n/2}/\Gamma(\frac n2)$
for large $n$, so the third term converges for all $s$ even though $k < k_0$ may be less than $R$. The terms of the
first series are $O(R^n u_0^{n/2}/\Gamma(n/2))$ away from the poles. Hence both sides are meromorphic in $\CC$
and agree for $\operatorname{Re}s > 1$. The residue of the first term of $M$ at $s = 1-n$ is
$a_n\sqrt\pi/\Gamma(\frac12) = a_n$.
\end{proof}

\begin{example}[Riemann]\label{ex:riemann}
For $f = 1$, $k_0 = 1$ and $u_0 = \pi$, the substitution $u = \pi^2m^2/v$ gives
$\sqrt\pi J_m(s) = \pi^{s-1/2}m^{s-1}\Gamma(\frac{1-s}2,\pi m^2)/\Gamma(\frac s2)$, and
\begin{multline*}
  \zeta(s) = \sum_{k=1}^{\infty}k^{-s}Q\Bigl(\frac s2,\pi k^2\Bigr)
  + \frac{\pi^{s-1/2}\Gamma(\frac{1-s}2)}{\Gamma(\frac s2)}\sum_{m=1}^{\infty}m^{s-1}Q\Bigl(\frac{1-s}2,\pi m^2\Bigr) \\
  + \frac{\pi^{s/2}}{\Gamma(\frac s2)\,s(s-1)},
\end{multline*}
which is Riemann's formula \cite{Rie59}; the symmetry of the two sums under $s \mapsto 1-s$ is the functional equation.
\end{example}

\begin{remark}[Poisson summation]\label{rem:poisson}
Theorem~\ref{thm:split} is a smoothed form of the Poisson summation formula
\[
  \sum_{k\in\ZZ}G_s(k) = \sum_{m\in\ZZ}\widehat G_s(m), \qquad G_s(x) = |x|^{-s}F(|x|).
\]
Formally, for $u_0 = \infty$ the dual terms become
\[
  \sqrt\pi\int_0^\infty u^{-1/2}e^{-\pi^2m^2/u}h_s(u)\dd u = \widehat G_s(m);
\]
for $F = 1$ this is $\pi^{s-1/2}\Gamma(\frac{1-s}{2})\Gamma(\frac s2)^{-1}m^{s-1}$. The functional equation of $\zeta$ expresses that
$\widehat G_s$ is again a power. For a general $F$, $\widehat G_s$ is not of the form $|\xi|^{s'}F'(|\xi|)$ with $F'$
regular, and there is no functional equation; but the split point $u_0$ makes both sides converge rapidly.
\end{remark}

\subsection{The two sides are holonomic}
For fixed $s$, the $k$-terms are values $I_k(s) = \mathcal L_s(k^2)$ of the Laplace transform
$\mathcal L_s(K) = \int_{u_0}^\infty e^{-Ku}h_s(u)\dd u$, and the dual terms are values
$J_m(s) = \mathcal D_s(\pi^2m^2)$ of $\mathcal D_s(c) = \int_0^{u_0}u^{-1/2}e^{-c/u}h_s(u)\dd u$. With
\eqref{eq:kernelparts},
\[
  \mathcal D_s(c) = Z_{(1-s)/2}[g_0](c) + Z_{1-s/2}[g_1](c), \qquad
  Z_t[g](c) = \int_0^{u_0}u^{-t}e^{-c/u}g(u)\dd u.
\]

\begin{proposition}[{\cite[Lemma~3.5]{Joh26}}]\label{prop:dual}
If $g$ satisfies $\vartheta G = A(u)G + b(u)$ as in Lemma~\ref{lem:kernelhol}, then for every $t$ the moments
$Z_{t-i}[G](c)$, $0 \le i \le \deg A$, together with finitely many moments of powers of $u$ and the constants
$e^{-c/u_0}$ and $e^{-c/u_0}G(u_0)$, form a vector $Y(c)$ with $cY'(c) = (B_0 + cB_1)Y(c)$, where $B_0, B_1$ have
entries in $\QQ(s, u_0^{1/2})$. In particular $\mathcal D_s$ is holonomic, analytic in $\operatorname{Re}c > 0$,
with its only finite singular point at $c = 0$.
\end{proposition}

The proof is integration by parts, $cZ_{t+1}[g] = e^{-c/u_0}u_0^{1-t}g(u_0) - (1-t)Z_t[g] - Z_t[\vartheta g]$, and
$-\partial_cZ_t = Z_{t+1}$. Similarly $\mathcal L_s(K)$ is holonomic in $K$, and on the $k$-side the series in $n$
for $I_k(s)$ has terms satisfying a linear recurrence in $n$ with coefficients in $\QQ(s,k,u_0)$, since
$Q(a+1,x) = Q(a,x) + x^ae^{-x}/\Gamma(a+1)$.

Thus Theorem~\ref{thm:split} expresses $\zeta_F(s)$ as a sum of values of a holonomic function at the squares
$k^2$, plus a sum of values of another holonomic function at the points $\pi^2m^2$, plus an explicit term; with
$u_0$ fixed, $O(\sqrt p)$ terms of each sum suffice for $p$ bits.

\begin{remark}[Periodic weights and twists]\label{rem:twist}
For $L_F(s,\chi)$ with $\chi$ of period $q$, one splits $\chi$ into even and odd parts and applies Poisson summation on
each residue class; the dual terms become $\hat\chi(m)\,\mathcal D_s(\pi^2m^2/q^2)$ with the Gauss sums
$\hat\chi(m) = \sum_{a\bmod q}\chi(a)e^{2\pi iam/q}$ (and $u^{-1/2}$ replaced by $u^{-3/2}$ for odd $\chi$); see
\cite[Proposition~3.3]{Joh26}. For $\chi(k) = e^{2\pi i\alpha k}$ with real $\alpha$ the dual terms are the same moments
at $c = \pi^2(m-\alpha)^2$, $m \in \ZZ$.
\end{remark}

\begin{remark}[Two variables]
If the coefficients $a_n = a_n(t)$ depend polynomially (or analytically, with uniform bounds) on a parameter $t$,
all pieces of Theorem~\ref{thm:split} are analytic in $(s,t)$, and $\sum_{k=1}^{\infty}F_t(k)k^{-s}$ is meromorphic in $\CC^2$
with polar divisor contained in $\bigcup_n\{s = 1-n\}$ and residues $a_n(t)$. For the families
$\sum_{k=1}^{\infty}G(k)^tk^{-s}$ of Section~\ref{sec:defs} this gives a new proof of the meromorphic continuation of the
Dirichlet series of polynomials of Mahler \cite{Mah28} along with a fast evaluation scheme; see also
Section~\ref{sec:phi}.
\end{remark}

%======================================================================
\section{Large imaginary part}\label{sec:larget}

Throughout this section $F$ is analytic at infinity \eqref{eq:f}, $s = \sigma + it$ with $t > 0$, and
$\chi(s) = 2^s\pi^{s-1}\sin(\pi s/2)\Gamma(1-s)$ is the factor in $\zeta(s) = \chi(s)\zeta(1-s)$.

\subsection{A Hardy--Littlewood approximate functional equation}

\begin{proposition}\label{prop:hl}
Let $0 < \sigma + n_0 < 1$ be fixed, $x, y \ge C$ with $2\pi xy = t$ and $x \ge k_0$. Then
\begin{multline*}
  \zeta_F(s) = \sum_{1\le k\le x}F(k)k^{-s} + \chi(s)\sum_{1\le m\le y}F\Bigl(\frac{t}{2\pi m}\Bigr)m^{s-1} \\
  + O\bigl(x^{-\sigma-n_0} + t^{1/2-\sigma-n_0}y^{\sigma+n_0-1}\bigr).
\end{multline*}
\end{proposition}

\begin{proof}[Proof sketch]
As in the classical proof for $\zeta$ \cite[\S4.12]{Tit86}: split off the terms $k \le N$ with $N$ large,
express the tail by \eqref{eq:zetavalue} and the Euler--Maclaurin formula, and apply the van der Corput
transformation \cite[Lemma~4.10]{Tit86} to the sum of $F(k)k^{-\sigma}e^{-it\log k}$ over $x < k \le N$. The phase
$-t\log x + 2\pi mx$ is stationary at $x_m = t/(2\pi m)$, and the amplitude $F(x)x^{-\sigma}$ varies slowly
($F^{(j)}(x) \ll x^{-n_0-j}$), so the $m$-th dual term is the classical one, $\chi(s)m^{s-1}$, times $F(x_m)$, with
a relative error $O(1/t)$. \todo{write out the error terms in full.}
\end{proof}

For $F(k) \sim a_{n_0}k^{-n_0}$ the dual coefficients are $F(t/2\pi m) \approx a_{n_0}(2\pi m/t)^{n_0}$, and since
$\chi(s+n_0) \approx (t/2\pi)^{-n_0}\chi(s)$ the formula reduces to the classical one for $a_{n_0}\zeta(s+n_0)$ to
leading order; the higher coefficients of $f$ enter through the values $F(t/2\pi m)$ at the stationary points.
Table~\ref{tab:hl} compares the error of Proposition~\ref{prop:hl} for Khinchin's $Z$ (with $x = y$) with the
error of the classical formula for $\zeta(s+2)$; without the dual sum the error would be of the size of the
dual sum itself.

\begin{table}[t]
\centering\small
\begin{tabular}{rrrrrr}
\toprule
$t$ & $\sigma$ & $|Z(s)\log 2|$ & dual sum & error, $Z$ & error, $\zeta(s+2)$ \\
\midrule
$1\,000$ & $-1.5$ & $1.022$ & $0.657$ & $0.098$ & $0.114$ \\
$1\,000$ & $-1.2$ & $0.355$ & $0.187$ & $0.046$ & $0.053$ \\
$4\,000$ & $-1.5$ & $0.341$ & $1.338$ & $0.097$ & $0.104$ \\
$4\,000$ & $-1.2$ & $0.131$ & $0.409$ & $0.037$ & $0.039$ \\
$16\,000$ & $-1.5$ & $0.459$ & $0.547$ & $0.052$ & $0.054$ \\
$16\,000$ & $-1.2$ & $0.324$ & $0.137$ & $0.016$ & $0.017$ \\
\bottomrule
\end{tabular}
\caption{Proposition~\ref{prop:hl} for $F(k) = p_k\log 2$ ($n_0 = 2$) with $x = y = \sqrt{t/2\pi}$: absolute values of
the function, of the dual sum, and of the error, compared with the error of the classical approximate functional
equation for $\zeta(s+2)$ with the same $x, y$.}
\label{tab:hl}
\end{table}

\begin{remark}
Theorem~\ref{thm:split} with real $u_0$ is not efficient for large $t$: $\Gamma(\frac s2)^{-1}$ is of size
$e^{\pi t/4}$ and the terms cancel to that extent. For $L$-functions this is overcome by rotating the
integration path, i.e.\ taking $u_0$ complex \cite{Rub05}, which gives Riemann--Siegel-like formulas with
$O(\sqrt t)$ terms. Since all pieces of Theorem~\ref{thm:split} are holonomic in $u_0$, the same should work
here. \todo{check; this would give an exact, smoothed version of Proposition~\ref{prop:hl}.}
\end{remark}

\subsection{Zeros}\label{sec:zeros}

\begin{proposition}\label{prop:zerofree}
\begin{enumerate}
\item There is $\sigma_1$ such that $\zeta_F(s) \ne 0$ for $\operatorname{Re}s \ge \sigma_1$.
\item For every $\varepsilon > 0$ and $\sigma_2 < -n_0-\varepsilon$ there is $t_0$ such that $\zeta_F(s) \ne 0$ for
$\sigma_2 \le \operatorname{Re}s \le -n_0-\varepsilon$, $|\operatorname{Im}s| \ge t_0$.
\end{enumerate}
\end{proposition}

\begin{proof}
(1) If $k_1$ is the least $k$ with $F(k) \ne 0$, then $k_1^s\zeta_F(s) \to F(k_1)$ as $\operatorname{Re}s \to \infty$,
uniformly in $\operatorname{Im}s$. (2) In the strip, write $\zeta_F$ as in the proof of
Proposition~\ref{prop:growth}. By the functional equation and Stirling's formula,
$|\zeta(s+n_0)| \asymp t^{1/2-\sigma-n_0}|\zeta(1-s-n_0)| \gg t^{1/2+\varepsilon}$, since
$\operatorname{Re}(1-s-n_0) \ge 1+\varepsilon$. Each $\zeta(s+n)$ with $n > n_0$ is
$O(t^{1/2-\sigma-n_0-1}+t^{\varepsilon/2})$, and everything else is bounded, so $\zeta_F(s) \ne 0$ for large $t$.
\end{proof}

\begin{theorem}\label{thm:count}
For $0 < \varepsilon < 1$ let $N_F(T)$ be the number of zeros $\rho$ of $\zeta_F$, counted with multiplicity, with
$\operatorname{Re}\rho \ge -n_0-\varepsilon$ and $0 < \operatorname{Im}\rho \le T$. Then
\[
  N_F(T) = \frac{T}{2\pi}\log\frac{T}{2\pi e k_1} + O(\log T),
\]
where $k_1$ is the least $k$ with $F(k) \ne 0$. In particular, if $F(1) \ne 0$, then $N_F(T) - N(T) = O(\log T)$,
where $N(T)$ counts the zeros of $\zeta$.
\end{theorem}

\begin{proof}[Proof sketch]
By Proposition~\ref{prop:zerofree} the zeros counted lie, apart from finitely many, in the rectangle
$-n_0-\varepsilon \le \sigma \le \sigma_1$, $t_0 \le t \le T$. Apply the argument principle to its boundary. On the
right side the change of argument is $-T\log k_1 + O(1)$, since $\operatorname{Re}\bigl(k_1^s\zeta_F(s)/F(k_1)\bigr) > \frac12$
there. On the
left side, by the proof of Proposition~\ref{prop:zerofree}, $\zeta_F(s) = a_{n_0}\zeta(s+n_0)(1 + o(1))$, so the change
of argument is that of $\zeta(-\varepsilon+it)$ up to $O(1)$, which by the functional equation is that of
$\chi(-\varepsilon+it)$ up to $O(1)$, namely $-T\log(T/2\pi e) + O(1)$. On the top side the change of argument is
$O(\log T)$ by the usual Jensen argument \cite[\S9.4]{Tit86}, using the polynomial growth of $\zeta_F$
(Proposition~\ref{prop:growth}) and $|\zeta_F(\sigma_1+iT)| \gg 1$.
\end{proof}

(For $\zeta(s) - 1$, i.e.\ $k_1 = 2$, this is Landau's count of the $1$-points of $\zeta$.) So if $F(1) \ne 0$,
$\zeta_F$ has as many zeros as $\zeta$, up to $O(\log T)$. Where are they? On the line
$\operatorname{Re}s = \frac12 - n_0$ we have $\zeta_F(s) = a_{n_0}\zeta(s+n_0) + B(s)$, where
$B(s) = \zeta_F(s) - a_{n_0}\zeta(s+n_0,k_0) - \sum_{k=1}^{k_0-1}F(k)k^{-s}$ is an absolutely convergent Dirichlet series
for $\operatorname{Re}s > -n_0$, with coefficients $O(k^{-n_0-1})$. The zeros of $\zeta_F$ are thus the solutions of
$\zeta(w) = -B(w-n_0)/a_{n_0}$: ``$a$-points'' of $\zeta$ with a variable, bounded, almost periodic $a$. Levinson
\cite{Lev75} proved that for constant $a$ all but $O(N(T)/\log\log T)$ of the $a$-points with
$0 < \operatorname{Im}s \le T$ lie in $|\operatorname{Re}s - \frac12| < \varepsilon$. This suggests:

\begin{conjecture}\label{conj:cluster}
For every $\varepsilon > 0$, all but $o(T\log T)$ of the zeros of $\zeta_F$ with $0 < \operatorname{Im}\rho \le T$ and
$\operatorname{Re}\rho \ge -n_0 - \varepsilon$ satisfy $|\operatorname{Re}\rho - (\frac12 - n_0)| < \varepsilon$.
\end{conjecture}

There is no reason to expect a Riemann hypothesis: without an Euler product, and as for the Davenport--Heilbronn
examples \cite{DH36}, there should be $\gg T$ zeros off the line.

\subsection{Numerical experiments}
For Khinchin's $Z$ ($n_0 = 2$), a search in $-6 \le \sigma \le 3$, $0 < t \le 50$ finds exactly one zero near each
zero $\frac12 + i\gamma_n$ of $\zeta(s+2)$, i.e.\ near $-\frac32 + i\gamma_n$, but displaced to the right by about
$0.5$--$0.7$ (Table~\ref{tab:zeros}), and three zeros of small height. To first order, the zero near
$-\frac32 + i\gamma_n$ is displaced by $-B(-\frac32 + i\gamma_n)/\zeta'(\frac12+i\gamma_n)$, which is small only when
$|\zeta'|$ is large; at these heights the first-order prediction gives the direction but underestimates the shift.
Higher up the displacement decreases slowly, in line with Conjecture~\ref{conj:cluster}: near
$\gamma_{100} = 236.52$, $\gamma_{1000} = 1419.42$ and $\gamma_{3000} = 3533.33$ the zeros of $Z$ have real parts
$-1.203$, $-1.138$ and $-1.273$.

\begin{table}[t]
\centering\small
\begin{tabular}{rrr@{\qquad}rrr}
\toprule
$\gamma_n$ & $\operatorname{Re}\rho$ & $\operatorname{Im}\rho$ & $\gamma_n$ & $\operatorname{Re}\rho$ & $\operatorname{Im}\rho$ \\
\midrule
$14.1347$ & $-0.80962$ & $14.10628$ & $37.5862$ & $-0.97200$ & $37.62163$ \\
$21.0220$ & $-0.87839$ & $21.04380$ & $40.9187$ & $-0.97864$ & $40.96155$ \\
$25.0109$ & $-0.90838$ & $24.96595$ & $43.3271$ & $-0.99413$ & $43.22684$ \\
$30.4249$ & $-0.93768$ & $30.50982$ & $48.0052$ & $-1.00783$ & $48.16616$ \\
$32.9351$ & $-0.95140$ & $32.84059$ & $49.7738$ & $-1.00367$ & $49.65618$ \\
\bottomrule
\end{tabular}
\caption{The zeros $\rho$ of Khinchin's $Z(s)$ with $5 < \operatorname{Im}\rho \le 50$, next to the ordinates
$\gamma_n$ of the zeros of $\zeta$.}
\label{tab:zeros}
\end{table}

The three zeros of small height, $-1.5119 + 1.6018i$, $-3.0436 + 2.1068i$ and $-4.5285 + 2.4371i$, start a chain
continuing to the left, approximately $-5.9950 + 2.6727i$, $-7.4513 + 2.8396i$, $-8.8985 + 2.9510i$,
$-10.3363 + 3.0160i$, $-11.7643 + 3.0414i$, $-13.1822 + 3.0323i$, $-14.5900 + 2.9928i$, with real parts spaced by
about $1.45$ and bounded imaginary parts, between the poles at the negative integers. They play the role of the
trivial zeros of $\zeta(s+2)$. Heuristically, far to the left the first dual term dominates and the zeros approach
those of $s \mapsto \widehat G_s(1)$, a Mellin transform of $F(x)\cos(2\pi x)$; for $F = 1$ these are the zeros of
$\sin(\pi s/2)$. In particular $\zeta_F$ may have infinitely many zeros with bounded imaginary part, which is why
Theorem~\ref{thm:count} restricts to $\operatorname{Re}\rho \ge -n_0-\varepsilon$.
\todo{prove or disprove that the chain is infinite.}

%======================================================================
\section{Computation}\label{sec:complexity}

``Time'' means bit complexity on a multitape Turing machine; $\Mul(p)$ is a bound for the cost of multiplying
$p$-bit integers satisfying the usual regularity conditions ($\Mul(p)/p$ nondecreasing and
$\Mul(p\log p) = O(\Mul(p)\log p)$); one can take $\Mul(p) = p\log p$ \cite{HvdH21} or even
$p(\log p)^{1-\kappa}$ for an explicit $\kappa > 0$ \cite{OAI26}.

\begin{theorem}\label{thm:complexity}
Let $F$ be a regular sequence whose function part is defined over $\Qbar$ and whose initial values
$F(1), \dots, F(k_0-1)$ are algebraic, let $\chi$ be periodic with algebraic values, let $s_0 \in \Qbar$, and let
$j \ge 0$ be fixed. Then the $j$-th coefficient of the Laurent expansion of $L_F(s,\chi)$ at $s_0$ can be computed to
$p$ bits in $O(p^{1/2}\Mul(p)\log^2 p)$ bit operations.
\end{theorem}

\begin{proof}[Proof sketch]
By Lemma~\ref{lem:reduction} it suffices to treat finitely many $\zeta_{F'}^{(i)}(s_0+\sigma;k_0)$ with $F'$
analytic at infinity and $\sigma$ algebraic, so assume \eqref{eq:f} with algebraic $a_n$ and an algebraic point $s_0$
(we write $\chi = 1$; see Remark~\ref{rem:twist}). Take $u_0$ the square of a fixed rational number and apply
Theorem~\ref{thm:split}, differentiated $j$ times in $s$. By the bounds in the theorem, $O(\sqrt p)$ $k$-terms
and dual terms suffice. Each is computable in $O(\Mul(p)\log^2p)$:
the $k$-terms by binary splitting of the series in $n$, whose terms satisfy recurrences over a fixed number field,
starting from $Q(\frac{s_0}{2}, k^2u_0)$ and $Q(\frac{s_0+1}2, k^2u_0)$; the dual terms by the bit-burst algorithm
applied to the system of Proposition~\ref{prop:dual}, from starting values at a rational $c_0$ that are again
series with P-recursive terms, $e^{-c_0/u_0}$, and one incomplete gamma value per class of exponents modulo $\ZZ$;
the derivatives in $s$ by carrying jets, which only enlarges the systems by a factor $j+1$; and $M$ and its Laurent
coefficients are explicit. Incomplete gamma functions $\Gamma(a,y)$ at algebraic $a$ and rational $y = O(p)$, and
$\Gamma(a)$, cost $O(\Mul(p)\log^2 p)$ by binary splitting. For details of each step see \cite[Theorems~3.1, 3.10
and~3.12]{Joh26}, where the same analysis is carried out for $s_0 = 0$ and for shifted exponents.
\todo{the companion paper assumes $\operatorname{Re}(s_0+\sigma) > 1$ in Theorem~3.12; check that nothing changes
for the continued pieces, in particular $Q(a,x)$ with $\operatorname{Re}a \le 0$.}
\end{proof}

\begin{remark}
\begin{enumerate}
\item For $s_0$ given only by approximations the series can no longer be summed by binary splitting and the cost
is $p^{2+o(1)}$; the same is true if the coefficients of $f$ are not algebraic.
\item The bound is not uniform in $s_0$: the constants grow with the height of $s_0$, and, for real $u_0$,
exponentially with $|\operatorname{Im}s_0|$ (see the remark after Proposition~\ref{prop:hl}).
\item In floating-point evaluation of $M$, the functions $P(a,x)$ for $k < k_0$ must be computed directly from the
lower series and not as $1 - Q(a,x)$, since they are multiplied by $a_nk^{-n}$, which may be exponentially large.
\end{enumerate}
\end{remark}

%======================================================================
\section{Special values}\label{sec:special}

\subsection{The Khinchin zeta function}\label{sec:khinchinzeta}
Let $Z(s) = \sum_{k=1}^{\infty}p_kk^{-s}$ as in \eqref{eq:Z}. Here
\[
  f(x) = \frac{2\log(1+x) - \log(1+2x)}{\log 2}, \qquad a_n = \frac{(-1)^n(2^n - 2)}{n\log 2},
\]
with $\rho = \frac12$, so $k_0 = 3$, and $a_0 = a_1 = 0$, $n_0 = 2$. By Proposition~\ref{prop:cont}, $Z$ is
meromorphic with simple poles at $s = -1, -2, -3, \dots$ of residues $1/\log 2$, $-2/\log 2$, $\frac{7}{2\log2}, \dots$,
and is analytic in $\operatorname{Re}s > -1$.

\begin{proposition}\label{prop:Z}
\begin{enumerate}
\item $Z(0) = 1$, $Z'(0) = -\log K$, and $K_p = Z(-p)^{1/p}$ for $p < 1$, $p \ne 0$.
\item $Z(s) = \dfrac{1}{(s+1)\log2} - \dfrac{1}{\log2} + O(s+1)$ as $s \to -1$. Equivalently, the H\"older means
satisfy $\log 2\,(K_p)^p = (1-p)^{-1} - 1 + O(1-p)$ as $p \to 1^-$.
\end{enumerate}
\end{proposition}

\begin{proof}
(1) is the definition of $K$ and of the H\"older means
\[
  K_p = \lim_{n\to\infty}\Bigl(\frac1n\sum_{i=1}^{n}a_i^p\Bigr)^{1/p},
\]
which for almost all real numbers equal $(\sum_{k=1}^{\infty}k^pp_k)^{1/p}$ by the ergodic theorem \cite{Khi35}. (2) By summation by parts
$\sum_{k=1}^{N}k\,(2\log(k+1) - \log k - \log(k+2)) = \log(N+1) - N\log\frac{N+2}{N+1} = \log N - 1 + o(1)$,
so $\sum_{k=1}^{\infty}(kp_k - \frac{1}{k\log2})$ converges to $-(1+\gamma)/\log 2$. Since $Z(s) - \zeta(s+2)/\log 2$ is analytic at
$s = -1$ with value $\sum_{k=1}^{\infty}(kp_k - \frac1{k\log 2})$, and $\zeta(s+2) = 1/(s+1) + \gamma + O(s+1)$, the constant term
is $-1/\log 2$.
\end{proof}

So the value at $0$ and the constant term at the first pole are elementary (both by telescoping), while the
derivative at $0$ is Khinchin's constant and the value at $1$ is the reciprocal of the harmonic mean. By
Theorem~\ref{thm:complexity} all H\"older means $K_p$ with rational $p < 1$, not only the $K_{-n}$, can be computed
to $N$ bits in time $\widetilde O(N^{3/2})$. Some values, computed with the zeta-value expansion as an illustration:
\begin{align*}
  K_{1/2} &= 4.5330951149778199744\ldots, & K_{-1/2} &= 2.0548408683571442932\ldots, \\
  K_{1/3} &= 3.6217632159994293888\ldots, & K_{-1/3} &= 2.2119276176511428815\ldots, \\
  K_{-2} &= 1.4503403284956304060\ldots, & K_{-3} &= 1.3135070786879857468\ldots
\end{align*}
\todo{verify against the literature; Bailey, Borwein and Crandall \cite{BBC97} give the $K_{-n}$.}

\subsection{A two-variable function}\label{sec:phi}
Since $k(k+2)/(k+1)^2 = 1 - (k+1)^{-2}$, Khinchin's constant also sits in the family
\[
  \Phi(s,t) = \sum_{k=1}^{\infty}k^{-s}\bigl(1 - (k+1)^{-2}\bigr)^t.
\]
Here $f_t(x) = ((1+2x)/(1+x)^2)^t$ is algebraic, analytic for $|x| < \frac12$, with coefficients
$a_0(t) = 1$, $a_1(t) = 0$, $a_2(t) = -t$, $a_3(t) = 2t$, \dots\ polynomial in $t$. By Theorem~\ref{thm:split}
(Remark on two variables), $\Phi$ is meromorphic in $\CC^2$ with polar divisor
$\{s=1\}\cup\bigcup_{n\ge2}\{s = 1-n\}$; there is no pole along $s = 0$. Moreover $\Phi(s,0) = \zeta(s)$,
$-\partial_t\Phi(s,0) = Z(s)\log 2$, and
\[
  \gamma = \CT_{s=1}\Phi(s,0), \qquad \log K\log 2 = \partial_s\partial_t\Phi(0,0).
\]
So Euler's and Khinchin's constants are coefficients of the same function at the points $(1,0)$ and $(0,0)$, and
the splitting of $\Phi$ is a one-parameter deformation of Riemann's formula (Example~\ref{ex:riemann}), to which it
reduces at $t = 0$.

\subsection{Further examples}
\begin{enumerate}
\item The rational examples
\[
  \sum_{k=2}^{\infty}\frac{k^{-s}}{k^2-1}, \qquad \sum_{k=1}^{\infty}\frac{k^{-s}}{k(k+1)}
\]
have derivatives $-\mathcal A$ and $-\mathcal C$ at $s = 0$; their poles are at $s = -1, -3, -5, \dots$ and at
$s = -1, -2, -3, \dots$. Partial fractions do not reduce the derivatives to values of $\zeta'$, since $\log k$ and
the denominator $k+1$ do not match.
\item The Hardy--Littlewood function: $F(k) = \sin(x/k)$ is entire in $1/k$, $k_0 = 1$, and
\[
  \sum_{k=1}^{\infty}\sin\Bigl(\frac xk\Bigr)k^{-s} = \sum_{m=0}^{\infty}\frac{(-1)^mx^{2m+1}}{(2m+1)!}\,\zeta(s+2m+1)
\]
is meromorphic with poles at $s = 0, -2, -4, \dots$, the pole at $s = -2m$ having residue
$(-1)^mx^{2m+1}/(2m+1)!$; $H(x)$ is its value at $s = 1$.
\item The vortex sum
\[
  \sum_{k=1}^{\infty} k^{-s}\,{}_2F_1\Bigl(\frac32,\frac32;2;-\frac{4}{\lambda^2k^2}\Bigr),
\]
with $\rho = \lambda/2$: the
velocity of a ring in a periodic array of coaxial vortex rings \cite{MS22} is $\lambda^{-3}$ times its value at
$s = 3$.
\item Planck--Larkin: $\sigma(T) = \zeta_F(-2)$ with $F(k) = e^{x/k^2} - 1 - x/k^2$; $f$ is entire, and the poles
are at $s = -3, -5, \dots$; $s = -2$ is a regular point, and $\sigma(T) = \sum_{k=1}^{\infty}k^2F(k)$ converges.
\end{enumerate}

\subsection{Computations}
The companion paper \cite{Joh26} reports computations with Theorem~\ref{thm:split} at $s = 0$ (and its derivative):
Khinchin's constant to $10^7$ digits (in $11.4$ hours on $8$ threads, against $65$ hours for the zeta-value
expansion), and $\mathcal A$, $\mathcal C$, $K_{-1}$, $U(1)$ and $H(1)$ to $10^6$ digits with a generic
implementation. \todo{decide how much of this to repeat here; add computations at other $s$, e.g.\ $K_{1/2}$ to
$10^6$ digits.}

%======================================================================
\section{Open problems}\label{sec:open}

\begin{enumerate}
\item \emph{Functional equations.} For which $F$ is $\widehat G_s$ again of the form $|\xi|^{s'}F'(|\xi|)$ with
$F'$ regular, giving a functional equation relating $\zeta_F$ and $\zeta_{F'}$? Besides powers, the inhomogeneous
Epstein (Chowla--Selberg) sums $\sum_{k=1}^{\infty}(k^2+a^2)^{-s}$, with modified Bessel functions on the dual side
\cite{CS67,BGMWZ13}, are of this kind, but only as diagonals of two-variable families.
\item \emph{Zeros.} Prove Conjecture~\ref{conj:cluster}, and determine whether a positive proportion of the zeros
lie off the line $\operatorname{Re}s = \frac12-n_0$. Is the chain of zeros of small height infinite?
\item \emph{Value distribution.} Is $\zeta_F$ universal (in the sense of Voronin) in the strip
$\frac12 - n_0 < \operatorname{Re}s < 1 - n_0$? Is it hypertranscendental?
\item \emph{Large height.} Give an exact, smoothed Riemann--Siegel formula for $\zeta_F$ with $O(\sqrt t)$ terms, and
an algorithm uniform in $\operatorname{Im}s$.
\item \emph{Closed forms.} Are there identities relating $Z'(0)$, i.e.\ $\log K$, to other periods or special
values? The derivatives $Z^{(j)}(0)$, $j \ge 1$, are the moments of $\log a_1$ under the Gauss measure.
\end{enumerate}

%======================================================================
% CHECK bibliographic details of DH36, Kat97, Ost20, Zag06, HL23 (added from memory/web search).
\begin{thebibliography}{99}
\bibitem{AV84} T.~M.~Apostol and T.~H.~Vu, Dirichlet series related to the Riemann zeta function, J.\ Number Theory 19 (1984), 85--102.
\bibitem{BBC97} D.~H.~Bailey, J.~M.~Borwein and R.~E.~Crandall, On the Khintchine constant, Math.\ Comp.\ 66 (1997), 417--431.
\bibitem{BGMWZ13} J.~M.~Borwein, M.~L.~Glasser, R.~C.~McPhedran, J.~G.~Wan and I.~J.~Zucker, \emph{Lattice Sums Then and Now}, Cambridge University Press, 2013.
\bibitem{CN81} P.~Cassou-Nogu\`es, Applications arithm\'etiques de l'\'etude des valeurs aux entiers n\'egatifs des s\'eries de Dirichlet associ\'ees \`a un polyn\^ome, Ann.\ Inst.\ Fourier 31 (1981), no.~4.
\bibitem{CS67} S.~Chowla and A.~Selberg, On Epstein's zeta-function, J.\ Reine Angew.\ Math.\ 227 (1967), 86--110.
\bibitem{Cra12} R.~E.~Crandall, Unified algorithms for polylogarithm, $L$-series, and zeta variants, preprint, 2012.
\bibitem{DH36} H.~Davenport and H.~Heilbronn, On the zeros of certain Dirichlet series, J.\ London Math.\ Soc.\ 11 (1936), 181--185 and 307--312.
\bibitem{Dok04} T.~Dokchitser, Computing special values of motivic $L$-functions, Experiment.\ Math.\ 13 (2004), 137--149.
\bibitem{Eli95} E.~Elizalde, \emph{Ten Physical Applications of Spectral Zeta Functions}, Lecture Notes in Physics m35, Springer, 1995.
\bibitem{Ess97} D.~Essouabri, Singularit\'es des s\'eries de Dirichlet associ\'ees \`a des polyn\^omes de plusieurs variables et applications en th\'eorie analytique des nombres, Ann.\ Inst.\ Fourier 47 (1997), 429--483.
\bibitem{FGD95} P.~Flajolet, X.~Gourdon and P.~Dumas, Mellin transforms and asymptotics: harmonic sums, Theoret.\ Comput.\ Sci.\ 144 (1995), 3--58.
\bibitem{HL23} G.~H.~Hardy and J.~E.~Littlewood, The approximate functional equation in the theory of the zeta-function, with applications to the divisor-problems of Dirichlet and Piltz, Proc.\ London Math.\ Soc.\ (2) 21 (1923), 39--74.
\bibitem{HvdH21} D.~Harvey and J.~van der Hoeven, Integer multiplication in time $O(n\log n)$, Ann.\ of Math.\ 193 (2021), 563--617.
\bibitem{Inc26} E.~L.~Ince, \emph{Ordinary Differential Equations}, Longmans, Green and Co., 1926.
\bibitem{Joh21} F.~Johansson, Rapid computation of special values of Dirichlet $L$-functions, arXiv:2110.10583, 2021.
\bibitem{Joh26} F.~Johansson, Computing infinite sums of holonomic functions in subquadratic time, in preparation, 2026.
\bibitem{Kat97} M.~Katsurada, An application of Mellin--Barnes' type integrals to the mean square of Lerch zeta-functions, Collect.\ Math.\ 48 (1997), 137--153.
\bibitem{Khi35} A.~Khintchine, Metrische Kettenbruchprobleme, Compositio Math.\ 1 (1935), 361--382.
\bibitem{Kir02} K.~Kirsten, \emph{Spectral Functions in Mathematics and Physics}, Chapman \& Hall/CRC, 2002.
\bibitem{Lav68} A.~F.~Lavrik, Approximate functional equations of Dirichlet functions, Izv.\ Akad.\ Nauk SSSR Ser.\ Mat.\ 32 (1968), 134--185.
\bibitem{Lev75} N.~Levinson, Almost all roots of $\zeta(s) = a$ are arbitrarily close to $\sigma = 1/2$, Proc.\ Nat.\ Acad.\ Sci.\ USA 72 (1975), 1322--1324.
\bibitem{Mah28} K.~Mahler, \"Uber einen Satz von Mellin, Math.\ Ann.\ 100 (1928).
\bibitem{MS22} E.~Masroor and M.~A.~Stremler, On the topology of the atmosphere advected by a periodic array of axisymmetric thin-cored vortex rings, Regul.\ Chaotic Dyn.\ 27 (2022), 183--197.
\bibitem{OAI26} OpenAI, Integer multiplication below $n\log n$, OpenAI Math Release preprint, September 2026, \url{https://github.com/openai/math/tree/main/preprints/Integer-multiplication-below-n-log-n-September-23-2026}.
\bibitem{Ost20} A.~Ostrowski, \"Uber Dirichletsche Reihen und algebraische Differentialgleichungen, Math.\ Z.\ 8 (1920), 241--298.
\bibitem{Pet98} M.~Peter, Dirichlet series associated with polynomials, Acta Arith.\ 84 (1998), 245--278.
\bibitem{Rie59} B.~Riemann, Ueber die Anzahl der Primzahlen unter einer gegebenen Gr\"osse, Monatsber.\ Berliner Akad., 1859.
\bibitem{Rub05} M.~Rubinstein, Computational methods and experiments in analytic number theory, in \emph{Recent Perspectives in Random Matrix Theory and Number Theory}, LMS Lecture Note Ser.\ 322, Cambridge University Press, 2005, 425--506.
\bibitem{Sar84} P.~Sargos, Prolongement m\'eromorphe des s\'eries de Dirichlet associ\'ees \`a des fractions rationnelles de plusieurs variables, Ann.\ Inst.\ Fourier 34 (1984), no.~3, 83--123.
\bibitem{Ste07} J.~Steuding, \emph{Value-Distribution of $L$-Functions}, Lecture Notes in Math.\ 1877, Springer, 2007.
\bibitem{Tit86} E.~C.~Titchmarsh, \emph{The Theory of the Riemann Zeta-Function}, 2nd ed., revised by D.~R.~Heath-Brown, Oxford University Press, 1986.
\bibitem{Zag06} D.~Zagier, The Mellin transform and related analytic techniques, appendix to E.~Zeidler, \emph{Quantum Field Theory I}, Springer, 2006, 305--323.
\end{thebibliography}

\end{document}
